\section{Spherical inert base change on an original member}

% BEGIN B85U1
\paragraph{B85U1. Reading the original unitary and cuspidal labels.}
Let $D$ be a selected member. Its norm-choice realization
$c.a.\mathrm{nt}$ contains, before the localization is added, the
Weil--Deligne identification $T$, the embedding of its cuspidal
labels in the unitary cuspidal spectrum (with its image condition),
and the realization of its unitary labels in automorphic functions.
They survive the localization: identify the underlying packaged
datum and character frame by the two projections of their equality,
then the representative by its equality in the member. The three
maps are $A.\mathrm{wdT}$, the same cuspidal embedding $\mathrm{cg}$,
and $A.\mathrm{au}$ with its realization $\mathrm{au}$.
These are the fields of the original member used in what follows.

For every permitted $j$ take any $\pi\in D.\mathrm{nbc}.\mathrm{AU}(j)$.
Its automorphic realization $U$ is an irreducible space of functions
on the definite unitary group in the base-change datum. Explicitly,
for a finite vector-valued form $f$ of weight $W$, take the span
of the scalar functions $\lambda(R(g_\infty)^{-1}f(g_f))$ with
$\lambda\in W^\vee$ and the algebraic realization $R$. Rational
left invariance is the transformation formula of $f$. Right finite
translations act on $f$ and right infinite translations act on the
dual coefficient. Evaluation detects a nonzero $f$; the map from the
irreducible tensor of the finite representation and the dual
archimedean factor therefore embeds it. Here irreducibility follows
by separating the matrix entries on the irreducible finite-dimensional
archimedean factor. Finite Hecke sums and local intertwiners are
transported by this embedding. The
realization of the split local component and the action along the
split-place isomorphism give isomorphic irreducible representations:
an intertwiner identifies their isomorphism classes, and the
injective Weil--Deligne identification gives the original local
label at this $w$. At a nonsplit $w$, with place $v$ below it,
the native spherical condition is carried by $A.\mathrm{sph}$ to
the predicate $\mathrm{SphAtG}$ on this realization. It includes
sphericity and, at a nonsplit place, goodness of its split local
supports. This second conclusion uses the definition of the member's
predicate. It is needed since only cuspidal representations with
the permitted supercuspidal support are enumerated. Choose local
conjugation and a uniformizer fixed by it from the unramified
quadratic completion. The hyperspecial double cosets then give
the same spherical Hecke eigencharacter in the row model and in
$U$, by the defining spherical Hecke action on its realized
functions. The normalized Satake readings agree up to the relative
Weyl group; in particular their image under unramified base change
agrees. The algebraic weight reading also agrees, by transport of
the representation at infinity and its contragredient weight pairing.
All three comparisons concern the same $U$.

On the general linear side, a normalized conjugate-self-dual
cuspidal representation with allowable support at every finite
place defines a unitary cuspidal class: the normalized discrete
cuspidal realization is unitary. This class has a label by the
image condition of $\mathrm{cg}$. Its local reading is the actual
local Langlands parameter, by uniqueness of local components in
the cuspidal representation and the local clause of $\mathrm{cg}$. A72F7 identifies the allowed
weak local correspondence with classical local Langlands on exactly
these supports.
The infinity and conjugate-self-duality clauses give the ordinary
infinitesimal character and the same global duality for the label.
% END B85U1

% BEGIN B85U2
\paragraph{B85U2. The original spherical inert base-change precondition.}
Assume the native field hypotheses of $\mathrm{Conclusion1}$ for
spherical inert data. Fix a permitted definite row of positive rank
$n$ and an original unitary label spherical at every nonsplit place.
B85U1 gives its definite automorphic realization $U$. Chebotarev for
the quadratic CM extension gives an inert place; there the
sphericity hypothesis supplies the split-support conclusion of
B85U1 for every split place. Apply A72F6 (strong definite unitary
base change, including the ascending Speh convention) to $U$.
It gives discrete conjugate-self-dual general linear constituents
whose ranks sum to $n$, with the local transfer at split and
spherical inert places and the specified infinity character.
Write their discrete-spectrum parameters as $(P_i,d_i)$, with
$P_i$ unitary cuspidal of rank $a_i>0$, $d_i>0$. Thus the rank
of the $i$th part is $a_i d_i$.

Each $P_i$ has a label in the original cuspidal enumeration.
Indeed at a split place the cuspidal support of the normalized
product is the union, with multiplicities, of the supports of the
local constituents. Its good support from B85U1 therefore passes
to each of them. At a nonsplit place the transferred representation
is spherical, whose supercuspidal support consists of unramified
characters of rank one. In either case passage from the Speh
constituent to its cuspidal base $P_i$ preserves goodness: the
support of the Speh quotient is the union of the supports of the
$d_i$ unramified translates of its cuspidal base. The predicate
on an individual support element is stable under these twists by
its definition (rank one, or a smooth character twist of a
principal-level-fixed representation at the specified local field).
Local uniqueness identifies the chosen base at every place.
Consequently B85U1 gives a single original cuspidal label $\pi_i$
with those local and archimedean readings.

The original list consists of the triples $(a_i,\pi_i,d_i)$.
The local Langlands multiset of a Speh quotient is the direct sum
of the parameters with exponents $(d_i-1)/2,(d_i-3)/2,\ldots,
(1-d_i)/2$. For clarity, the ascending discrete embedding gives the descending
quotient by local duality and contragredience. By generic unitary
classification the square-integrable segments of the cuspidal
local factor have real exponents strictly between $-1/2$ and
$1/2$. Translates spaced one apart therefore concatenate in
descending standard order. The unique irreducible head has the
union of those translated segment parameters, including their
nilpotent operators. The parameter of the unitary product
is the sum of these parameters by normalized induction and the
unitary product irreducibility of V72L1 and the standard quotients of V72L3.
The Satake normalization on a spherical representation is the
unramified Langlands parameter. Thus at split places the sum is
the given local label of $U$. At a nonsplit place it is unramified
with eigenvalues $\alpha_j,\alpha_j^{-1}$ and an additional $1$
in odd rank, the prescribed $\mathrm{ubc}$ of the member's
spherical tuple. Geometric reciprocity inverts the Satake roots;
inversion preserves this multiset, with multiplicities. At infinity the same translates of the infinitesimal
characters sum to the character prescribed by the unitary weight.
This last multiset has no repetitions, since a weakly dominant
weight becomes strictly ordered after adding the half sum of
positive roots. Two equal pairs $((a_i,\pi_i),d_i)$ would supply
the same nonempty submultiset twice there. This proves the required
distinctness as well. All clauses now hold for this one list for
every original $j,\pi$ in the test. This proves the
$\texttt{bc}$ field of the selected Galois preconditions.
% END B85U2
